% ============================================================
% PRESA D'ARIA TBCC - VERSIONE BIANCO E NERO
% ============================================================


% ============================================================
% 1. PARAMETRI PRINCIPALI
% ============================================================

% Dati reali MATLAB
\def\hCatt{1.1005}
\def\LSuper{2.2532}

\def\LRampaOneTwo{1.4907}
\def\LRampaTwoThree{0.7145}
\def\LRampaThreeFour{0.4435}

% Delta s
\def\DeltaS{0.2253}

% Angoli MATLAB originali
% Servono SOLO per costruire la geometria.
% Non vengono visualizzati.
\def\deltaOneRaw{0.2057}
\def\deltaTwoRaw{0.2521}
\def\deltaThreeRaw{0.3093}

% Mach
\def\Mzero{3.5}

% Aree
\def\AtwoRam{0.5273}

% Sostituire con motore.presaAria.A2 se hai il valore preciso
\def\AtwoTurbo{0.5273}

% Scala globale
\def\UnitScale{2.45}



% ============================================================
% 2. POSIZIONI MODIFICABILI DELLE SCRITTE
% ============================================================
%
% DX positivo = destra
% DX negativo = sinistra
%
% DY positivo = alto
% DY negativo = basso
% ============================================================

% Shock point
\def\ShockTextDX{-0.20}
\def\ShockTextDY{-0.08}

% Mach
\def\MachTextDX{-0.10}
\def\MachTextDY{0.00}

% Compressore
\def\CompressorTextDX{-0.15}
\def\CompressorTextDY{0.00}

% Post combustore
\def\PostTextDX{0.00}
\def\PostTextDY{0.00}

% B
\def\BTextDX{-0.04}
\def\BTextDY{0.06}

% C
\def\CTextDX{0.00}
\def\CTextDY{0.07}

% D
\def\DTextDX{0.00}
\def\DTextDY{0.07}

% E
\def\ETextDX{0.00}
\def\ETextDY{0.07}

% A = Shock point
\def\ATextDX{-0.08}
\def\ATextDY{0.04}

% A' = punto dopo Delta_s
\def\AprimeTextDX{0.05}
\def\AprimeTextDY{0.04}



% ============================================================
% 3. CONVERSIONE ANGOLI
% ============================================================

% I valori MATLAB vengono usati con sin/cos,
% quindi li convertiamo internamente per TikZ.

\pgfmathsetmacro{\deltaOne}{deg(\deltaOneRaw)}
\pgfmathsetmacro{\deltaTwo}{deg(\deltaTwoRaw)}
\pgfmathsetmacro{\deltaThree}{deg(\deltaThreeRaw)}

\pgfmathsetmacro{\thetaFlow}{
    \deltaOne+\deltaTwo+\deltaThree
}



% ============================================================
% 4. ALTEZZA COMPRESSORE
% ============================================================

\pgfmathsetmacro{\hComp}{
    2*sqrt(\AtwoTurbo/pi)
}



% ============================================================
% 5. COORDINATE RAMPE
% ============================================================

% B
\pgfmathsetmacro{\xB}{0}
\pgfmathsetmacro{\yB}{\hCatt}

% C
\pgfmathsetmacro{\xC}{
    \xB+\LRampaOneTwo*cos(\deltaOne)
}

\pgfmathsetmacro{\yC}{
    \yB-\LRampaOneTwo*sin(\deltaOne)
}

% D
\pgfmathsetmacro{\xD}{
    \xC+\LRampaTwoThree*cos(\deltaOne+\deltaTwo)
}

\pgfmathsetmacro{\yD}{
    \yC-\LRampaTwoThree*sin(\deltaOne+\deltaTwo)
}

% E
\pgfmathsetmacro{\xE}{
    \xD+\LRampaThreeFour*cos(\thetaFlow)
}

\pgfmathsetmacro{\yE}{
    \yD-\LRampaThreeFour*sin(\thetaFlow)
}



% ============================================================
% 6. SHOCK POINT = A
% ============================================================

\pgfmathsetmacro{\xShock}{\LSuper}
\pgfmathsetmacro{\yShock}{0}



% ============================================================
% 7. PUNTO A' DOPO DELTA S
% ============================================================

% Delta_s è interpretato come distanza orizzontale,
% come nella figura di riferimento.

\pgfmathsetmacro{\xLip}{
    \xShock+\DeltaS
}

% Il punto continua nella direzione della terza rampa
\pgfmathsetmacro{\yLip}{
    \yShock-\DeltaS*tan(\thetaFlow)
}



% ============================================================
% 8. GEOMETRIA INTERNA
% ============================================================

\pgfmathsetmacro{\xHinge}{
    \xE+0.80*\LSuper
}

\pgfmathsetmacro{\xMax}{
    \xHinge+0.25*\LSuper
}

\pgfmathsetmacro{\yFloor}{
    -0.45*\hCatt
}

\pgfmathsetmacro{\yHinge}{
    \yFloor+sqrt(\AtwoRam/pi)
}

\pgfmathsetmacro{\yRoof}{
    \yHinge+\hComp
}



% ============================================================
% 9. BEZIER - TETTO
% ============================================================

\pgfmathsetmacro{\RoofPoneX}{
    \xE+0.5*(\xHinge-\xE)
}

\pgfmathsetmacro{\RoofPoneY}{
    \yE
}

\pgfmathsetmacro{\RoofPtwoX}{
    \xHinge-0.2*(\xHinge-\xE)
}

\pgfmathsetmacro{\RoofPtwoY}{
    \yRoof
}



% ============================================================
% 10. BEZIER - PARATIA
% ============================================================

\pgfmathsetmacro{\dxFlap}{
    \xHinge-\xE
}

\pgfmathsetmacro{\FlapPoneX}{
    \xE+0.35*\dxFlap
}

\pgfmathsetmacro{\FlapPoneY}{
    \yE-0.25*\dxFlap*tan(\thetaFlow)
}

\pgfmathsetmacro{\FlapPtwoX}{
    \xHinge-0.35*\dxFlap
}

\pgfmathsetmacro{\FlapPtwoY}{
    \yHinge
}



% ============================================================
% 11. BEZIER - PAVIMENTO
% ============================================================

\pgfmathsetmacro{\dxFloor}{
    \xHinge-\xLip
}

\pgfmathsetmacro{\FloorPoneX}{
    \xLip+0.35*\dxFloor
}

\pgfmathsetmacro{\FloorPoneY}{
    \yLip-0.25*\dxFloor*tan(\thetaFlow)
}

\pgfmathsetmacro{\FloorPtwoX}{
    \xHinge-0.35*\dxFloor
}

\pgfmathsetmacro{\FloorPtwoY}{
    \yFloor
}



% ============================================================
% 12. FIGURA
% ============================================================


\centering

\begin{tikzpicture}[
    x=\UnitScale cm,
    y=\UnitScale cm,
    line cap=round,
    line join=round,
    font=\scriptsize,
    >={Stealth[length=2.5mm,width=1.8mm]}
]

\clip
  ({-0.42*\LSuper},{\yFloor-0.45})
  rectangle
  ({\xHinge-0.65*\LSuper},{\hCatt+0.15});

% ============================================================
% FUSOLIERA SUPERIORE / REFERENCE AXIS
% ============================================================

\draw[
    black,
    line width=0.9pt
]
({-0.15*\LSuper},\hCatt)
--
(\xMax,\hCatt);





% ============================================================
% RAMPE B-C-D-E
% ============================================================

\draw[
    black,
    line width=1.3pt
]
(\xB,\yB)
--
(\xC,\yC)
--
(\xD,\yD)
--
(\xE,\yE);



% ============================================================
% ONDE D'URTO - TUTTE NERE
% ============================================================

\draw[
    black,
    line width=1.1pt
]
(\xB,\yB)
--
(\xShock,\yShock);


\draw[
    black,
    line width=1.1pt
]
(\xC,\yC)
--
(\xShock,\yShock);


\draw[
    black,
    line width=1.1pt
]
(\xD,\yD)
--
(\xShock,\yShock);


\draw[
    black,
    line width=1.1pt
]
(\xE,\yE)
--
(\xShock,\yShock);



% ============================================================
% SHOCK POINT
% ============================================================




% ============================================================
% TRATTO A -> A'
% ============================================================

\draw[
    black,
    dashed,
    line width=1pt
]
(\xShock,\yShock)
--
(\xLip,\yLip);



% ============================================================
% ETICHETTE A E A'
% ============================================================

\node
at
(
    {\xShock+\ATextDX},
    {\yShock+\ATextDY-0.1}
)
{
    $A$
};


\node
at
(
    {\xLip+\AprimeTextDX},
    {\yLip+\AprimeTextDY}
)
{
    $A'$
};



% ============================================================
% TETTO
% ============================================================

\draw[
    black,
    line width=1.3pt
]
(\xE,\yE)
..
controls
    (\RoofPoneX,\RoofPoneY)
and
    (\RoofPtwoX,\RoofPtwoY)
..
(\xHinge,\yRoof)
--
(\xMax,\yRoof);



% ============================================================
% PARATIA MOBILE
% ============================================================

\draw[
    black,
    line width=1.6pt
]
(\xE,\yE)
..
controls
    (\FlapPoneX,\FlapPoneY)
and
    (\FlapPtwoX,\FlapPtwoY)
..
(\xHinge,\yHinge);


\fill[
    black
]
(\xHinge,\yHinge)
circle[radius=1.8pt];



% ============================================================
% PAVIMENTO
% ============================================================

\draw[
    black,
    line width=1.3pt
]
(\xLip,\yLip)
..
controls
    (\FloorPoneX,\FloorPoneY)
and
    (\FloorPtwoX,\FloorPtwoY)
..
(\xHinge,\yFloor)
--
(\xMax,\yFloor);



% ============================================================
% LINEA DIVISORIA
% ============================================================

\draw[
    black,
    line width=1.2pt
]
(\xHinge,\yHinge)
--
(\xMax,\yHinge);



% ============================================================
% FACCIA COMPRESSORE
% ============================================================

\draw[
    black,
    line width=2.3pt
]
(\xMax,\yHinge)
--
(\xMax,\yRoof);








% ============================================================
% ETICHETTE B C D E
% ============================================================

\node
at
(
    {\xB+\BTextDX},
    {\yB+\BTextDY}
)
{
    $B$
};


\node
at
(
    {\xC+\CTextDX},
    {\yC+\CTextDY}
)
{
    $C$
};


\node
at
(
    {\xD+\DTextDX},
    {\yD+\DTextDY}
)
{
    $D$
};


\node
at
(
    {\xE+\ETextDX},
    {\yE+\ETextDY}
)
{
    $E$
};



% ============================================================
% QUOTA VERTICALE h
% ============================================================

% Posizione della quota
\pgfmathsetmacro{\xHeightDim}{
    -0.38*\LSuper
}

% Linee di estensione
\draw[
    black,
    thin
]
(\xHeightDim,\hCatt)
--
(\xB,\hCatt);


\draw[
    black,
    thin
]
(\xHeightDim,\yShock)
--
(\xShock,\yShock);


% Freccia verticale
\draw[
    <->,
    black,
    line width=0.9pt
]
(\xHeightDim,\yShock)
--
(\xHeightDim,\hCatt);


% Etichetta h
\node[
    rotate=90,
    fill=white,
    inner sep=2pt
]
at
(
    \xHeightDim,
    {(\hCatt+\yShock)/2}
)
{
    $h$
};



% ============================================================
% QUOTA L_SUPER
% ============================================================

\pgfmathsetmacro{\yLsuperDim}{
    \yFloor-0.28
}

% Estensioni verticali
\draw[
    black,
    thin
]
(\xB,\yShock)
--
(\xB,\yLsuperDim);


\draw[
    black,
    thin
]
(\xShock,\yShock)
--
(\xShock,\yLsuperDim);


% Quota
\draw[
    <->,
    black,
    line width=0.9pt
]
(\xB,\yLsuperDim)
--
(\xShock,\yLsuperDim);


\node[
    fill=white,
    inner sep=2pt
]
at
(
    {(\xB+\xShock)/2},
    \yLsuperDim
)
{
    $L_{\mathrm{SUPER}}$
};



% ============================================================
% QUOTA DELTA_S
% ============================================================

\pgfmathsetmacro{\yDeltaDim}{
    \yLsuperDim+0.16
}

% Estensione shock point
\draw[
    black,
    thin
]
(\xShock,\yShock)
--
(\xShock,\yDeltaDim);


% Estensione A'
\draw[
    black,
    thin
]
(\xLip,\yLip)
--
(\xLip,\yDeltaDim);


% Quota Delta_s
\draw[
    <->,
    black,
    line width=0.9pt
]
(\xShock,\yDeltaDim)
--
(\xLip,\yDeltaDim);


\node[
    fill=white,
    inner sep=1pt
]
at
(
    {(\xShock+\xLip)/2},
    \yDeltaDim
)
{
    $\Delta_s$
};


\end{tikzpicture}


\caption{
Schema geometrico della presa d'aria TBCC.}
